\documentclass[12pt]{article}
\usepackage[utf8]{inputenc}
\usepackage{natbib}
\usepackage{amsmath,amssymb}
\usepackage{graphicx}

\title{Polarized Narratives and Overreactions in Economic Expectations: \ A Research Proposal}
\author{Lorenzo Gorini \\ Behavioral Macroeconomics}
\date{\today} 


\begin{document}

\maketitle

\begin{abstract}
    This proposal outlines a study on how \textbf{polarized narrative exposure} in media can lead to \textbf{diagnostic overreactions} in household economic expectations, and how these effects vary across socio-economic status (SES). I hypothesize that exposure to strongly partisan or polarized economic narratives causes households to form overly pessimistic or optimistic expectations (relative to fundamentals), consistent with the behavioral model of \emph{diagnostic expectations}. The research will exploit geographic variation in media exposure---treating narrative exposure as an intent-to-treat at the Designated Market Area (DMA) level---to identify causal effects on expectations. I will examine heterogeneous impacts by income, education, and other SES indicators to determine whether certain groups are more susceptible to narrative-driven expectation errors. The findings will contribute to our understanding of how media-driven narratives can amplify economic fluctuations through belief formation, and inform policies aimed at mitigating misinformation or communication gaps.
\end{abstract}

\section{Introduction and Motivation}
Economic expectations of households play a crucial role in macroeconomic outcomes, influencing consumption, investment, and policy effectiveness. Yet a growing body of evidence suggests that household expectations are not formed in a vacuum; they are influenced by the narratives and information to which people are exposed, particularly via mass media and social networks \citep{Carroll2003, Shiller2017}. In recent years, media environments have become increasingly polarized, with news outlets often presenting starkly different narratives about the economy depending on their political orientation \citep{Mullainathan2005, Martin2017}. This polarization raises the question: how does exposure to \emph{polarized economic narratives} affect the way people form expectations about economic variables such as inflation, unemployment, or growth?

Standard models of expectations in macroeconomics typically assume rational expectations or simple learning rules. However, these models struggle to explain certain patterns observed in survey data on household expectations. For instance, households often \emph{overreact} to salient news or recent experiences, swinging from optimism to pessimism more sharply than a fully rational agent would \citep{Bordalo2022}. Additionally, beliefs about the economy have been shown to diverge systematically along partisan or socio-economic lines, even when individuals face the same underlying economic reality \citep{Kamdar2023}. These observations hint at the role of \emph{narratives} and cognitive biases in expectation formation.

This research is motivated by an interplay of two key insights. First, narratives---coherent stories or frames that media use to explain economic events---can spread widely and influence behavior, as emphasized by the concept of \textit{Narrative Economics} \citep{Shiller2017}. When narratives become polarized (for example, one outlet consistently painting the economy as on the brink of recession while another touts robust growth), different populations are effectively living in different informational worlds. Second, a growing literature in behavioral economics documents \emph{diagnostic expectations} as a mechanism for expectation formation, wherein individuals overweight new information that is representative or salient, leading to forecast errors that are systematically biased in the direction of the news \citep{Bordalo2022}. This mechanism, based on the representativeness heuristic \citep{Tversky1974}, implies that people may \emph{overreact} to narratives that strongly signal a particular economic outcome.

The central question of this proposal is: \textbf{Does exposure to polarized economic narratives cause diagnostic overreactions in household expectations, and are these effects more pronounced for certain socio-economic groups?} Understanding this is important for several reasons. If partisan media narratives can cause segments of the population to become unduly optimistic or pessimistic about the economy, they could amplify business cycles or create self-fulfilling prophecies (e.g., widespread pessimism leading to reduced spending and a real downturn). Moreover, if these effects are heterogeneous---for instance, more severe for lower-income or less-educated households---it could imply that media-driven expectation errors have distributional consequences, potentially exacerbating economic inequality or unequal responses to shocks.

\section{Literature Review}
This project contributes to several strands of literature. First, it builds on research about media influence on economic and political outcomes. \citet{DellaVigna2007} provided early causal evidence that the introduction of a partisan news channel (Fox News) measurably shifted voting behavior, demonstrating that media can persuade and change beliefs. Subsequent work showed that media slant and availability can also polarize beliefs and attitudes: \citet{Martin2017}, for example, exploit quasi-random channel positioning in cable lineups as an instrument for exposure to Fox News and CNN, finding that slanted news affects not only voting patterns but also increases polarization in viewers' attitudes. These studies underscore that media content can causally impact what people think and do. However, most of this literature has focused on political outcomes; less is known about the effects on economic expectations and behavior. Second, this research draws on the emerging field of \textit{narrative economics} and the psychology of expectations. \citet{Shiller2017} argues that popular narratives can go "viral" and drive major economic events by influencing collective behavior. There is growing empirical evidence linking narratives to economic sentiment: for instance, \citet{Kamdar2023} document that households increasingly justify their economic outlook using partisan narratives, and that changes in political leadership (e.g., a new president) lead to mirror-opposite shifts in optimism between Republican- and Democrat-aligned households. Standard expectation models (like rational or purely extrapolative expectations) cannot fully explain these patterns of partisan divergence in sentiment. My study will add to this by explicitly testing whether narrative exposure itself (as opposed to just personal partisan identity) can cause overreaction in beliefs. Third, this project is grounded in behavioral macroeconomic theory, particularly the concept of \textit{diagnostic expectations} \citep{Bordalo2022}. Diagnostic expectations provide a micro-founded way to formalize overreaction: agents overweight the probability of outcomes that have become more salient or representative given recent news. In other words, if a particular narrative strongly suggests a recession, people applying the representativeness heuristic might assign an unduly high probability to a recession, more than a rational Bayesian would. Over time, as reality unfolds differently, these expectations are proven too extreme, leading to predictable reversals. This framework has been used to explain phenomena like credit cycles and asset price boom-bust dynamics, where investors get overly optimistic during booms and excessively pessimistic during busts \citep{Bordalo2022}. My contribution is to apply this idea to household macroeconomic expectations and to link it with heterogeneous media exposure. Finally, this study touches on the literature about heterogeneity in expectation formation. Research has shown that demographic factors and cognitive abilities can affect how people form and update beliefs. \citet{DAcunto2019} find that individuals with higher cognitive abilities have systematically lower inflation forecast errors, suggesting that information processing plays a role. Relatedly, there is evidence that less-educated or lower-income households may rely more on salient news (like grocery prices or headline news) when forecasting inflation, potentially making them more susceptible to media influence. Recent work by \citet{Polattimur2025} indeed finds that televised news coverage of inflation has a stronger effect on the inflation expectations of poorer and less-educated households, as well as younger viewers. This indicates that media-driven expectation biases may disproportionately impact certain socio-economic groups, possibly due to differences in information sources or financial literacy. My project will extend this line of inquiry by focusing on a broad set of economic expectations (beyond inflation) and explicitly linking heterogeneity to the \emph{content} of narratives (e.g., optimistic vs pessimistic slant) that different groups consume.

\section{Conceptual Framework}
The working hypothesis is that polarized narratives serve as shocks to people’s information sets, and due to diagnostic inference, individuals overweight these narrative signals. To illustrate, consider two regions or groups of viewers: one primarily exposed to an optimistic economic narrative (e.g., a news outlet that attributes all recent positive indicators to a booming economy), and another exposed to a pessimistic narrative (e.g., an outlet emphasizing downside risks and negative indicators). In a rational expectations framework, if both groups have the same underlying economic fundamentals, their forecasts for, say, next year’s GDP growth should be similar after processing all available information. However, under diagnostic expectations, each group will overreact in the direction of the narrative it sees. The optimistic-narrative group might forecast excessively high growth, while the pessimistic-narrative group forecasts excessively low or even negative growth, even if both have seen essentially the same aggregate data with different spins. Formally, one can think of an agent updating their belief $E_t[X_{t+1}]$ about a future outcome $X_{t+1}$ upon receiving a signal $s$. A rational Bayesian would weight this signal by its objective informativeness relative to prior beliefs. A diagnostically biased agent, in contrast, applies extra weight to signals that are representative of certain outcomes. If a narrative is very dramatic or aligns with a salient pattern (e.g., "this downturn looks like 2008 all over again"), the agent’s posterior belief overshoots. As \citet{Bordalo2022} put it, agents act "as if" the probability of extreme outcomes is higher than it truly is after seeing representative news. This mechanism can generate exaggerated swings in expectations across news cycles. Crucially, diagnostic overreaction is transient but impactful: beliefs will eventually correct as hard evidence accumulates, but in the interim they can induce real effects (like cutbacks in spending or investment due to unwarranted pessimism). This framework helps make sense of why surveys often find greater volatility in household expectations than can be justified by objective economic changes. It also differentiates our hypothesis from alternative models like rational inattention or sticky information, which often predict underreaction or delay in incorporating news \citep[e.g.,][]{Carroll2003}. Instead, here we posit an immediate but \emph{biased} reaction. \section{Empirical Strategy and Data}
To test the hypothesis, I will exploit geographic variation in media narrative exposure across the United States. The unit of analysis will often be at the Designated Market Area (DMA) level, a region that corresponds to a local media market (for example, the New York DMA or the Dallas-Fort Worth DMA). The key independent variable will be a measure of \textbf{polarized narrative exposure} for each DMA and time period. One approach to construct this measure is to leverage differences in cable news channel line-ups and viewership. Following \citet{Martin2017}, I can use the channel positions of news networks (like Fox News, CNN, MSNBC) on local cable systems as an instrument for viewership rates. The logic is that lower channel numbers (e.g., channel 10 vs channel 70) lead to higher casual viewership because viewers tend to start at low-number channels. These channel positions were historically set in ways unrelated to local political leanings, providing something close to a natural experiment in exposure. Thus, in some DMAs, Fox News ended up with a very favorable position (increasing exposure to its conservative narratives), while in others CNN might have the edge. I will combine this with data on the content of these networks, using text analysis of transcripts to quantify the economic sentiment or narrative slant on each network. For example, I will create an index of economic pessimism vs optimism in each network’s coverage (using keyword counts or machine learning methods on news transcripts). This content analysis will let me characterize each DMA’s “narrative exposure” at a given time—e.g., DMA A might have a high exposure to pessimistic narratives (if Fox is widely watched there during a recession scare), whereas DMA B has more optimistic narrative exposure (if CNN is emphasizing strong economy stories). Importantly, by treating the variation in exposure as happening at the DMA level (the intent-to-treat), I account for the fact that not every individual in a high-exposure area watches the news. The assumption is that enough people are influenced in high-exposure areas that the average expectations in that area shift, even if some individuals do not watch. This gives a conservative estimate of the narrative’s effect (an Intent-to-Treat effect), acknowledging that actual individual-level treatment is probabilistic. The dependent variables will be measures of household expectations. I plan to use survey data such as the University of Michigan’s Surveys of Consumers (which reports households’ expectations about inflation, unemployment, personal finances, etc.) or the Conference Board’s consumer confidence data, which can potentially be broken down by region. If microdata are available, I will geocode respondents to their DMA. Alternatively, if only aggregate indices by region are available, I will use those. I will also explore novel data sources: for instance, the New York Federal Reserve’s Survey of Consumer Expectations might be usable for certain breakdowns, or high-frequency polling data (like Gallup or Morning Consult) that include economic sentiment and can be segmented geographically. The empirical model will likely be a differences-in-differences or panel regression framework, exploiting changes in narrative exposure over time within a DMA. One potential design is to focus on specific events that caused a divergence in narratives. For example, major elections or policy announcements often trigger opposite spins on different networks. I can examine the pre- vs post-event change in expectations among respondents in areas more exposed to one network versus another. If, say, a shock (like an interest rate hike or a pandemic-related policy) is framed very negatively on one outlet and neutrally or positively on another, do we see a divergence in expectations (inflation expectations, economic outlook) between viewers in areas with higher exposure to the negative vs positive outlet? To address identification, I will use instrumental variables as described (channel positions as an instrument for actual narrative consumption) to mitigate self-selection bias (people choose what to watch) and reverse causality (local economic conditions influencing both media content and expectations). By controlling for DMA fixed effects and time fixed effects (and possibly DMA-specific linear trends), I aim to isolate the effect of changes in narrative exposure within a region over time. I will also control for actual economic conditions at the regional level (such as local unemployment rate or income growth) to ensure that any expectation differences are not simply driven by different fundamentals.

\textbf{Interpretation and Limitations.} The coefficient on narrative exposure (or its instrumented counterpart) in these regressions will be interpreted as the causal effect of narrative exposure on expectations, but it is important to remember this is an \emph{intent-to-treat} effect. Not everyone in a high-exposure DMA consumes the partisan narrative, so the true effect on those who do consume (the treatment-on-the-treated) is likely larger. I will discuss how to scale up to estimate the per-viewer effect, using known viewership rates. Another limitation is spillovers: people not directly watching might hear about the narratives second-hand (conversations, social media). In a sense, this means my ITT estimates might actually capture broader societal influence of narratives, which is relevant for macro outcomes, but it complicates a clean interpretation. I will attempt to mitigate concerns by checking, for example, if effects are stronger in demographic groups that are more likely to watch the news (e.g., the elderly) within the same DMA. If so, that strengthens the case that media is the channel.

\section{Heterogeneity by Socio-Economic Status}
A core part of the analysis is to examine how the effects of narrative exposure on expectations might differ by socio-economic status. There are several reasons to expect \textbf{SES heterogeneity}. Individuals with lower income or education may have fewer alternative information sources and might rely more heavily on television news for economic information, potentially making them more vulnerable to whatever narrative their preferred channel presents. They may also have lower baseline financial literacy or less time to seek out balanced information, which could exacerbate a representativeness heuristic in processing news. In contrast, higher-SES individuals might consume a broader diet of news (including financial news, newspapers, etc.), or may trust expert forecasts more, which could moderate narrative-induced swings. To test heterogeneity, I will interact the narrative exposure variable with indicators for education level (college educated vs not), income brackets, or occupation (e.g., blue-collar vs white-collar, if available). In survey microdata, this means running regressions at the individual level with DMA-level exposure as a key regressor and including interaction terms. For example, I will estimate models of the form:

\begin{align*}
    Expectation_{i,t} & = \alpha + \beta_1 \text{Exposure}_{r,t} + \beta_2 \text{LowSES}_i                                             \\
                      & + \beta_3 (\text{Exposure}_{r,t} \times \text{LowSES}_i) + \Gamma X_{it} + \Psi_r + \Lambda_t + \epsilon_{i,t}
\end{align*}

where $i$ indexes individuals, $r$ indexes region (DMA), and $t$ time. $LowSES_i$ might be a dummy for below-median income or not having a college degree. $\Psi_r$ and $\Lambda_t$ are region and time fixed effects. Here $\beta_1$ captures the effect of exposure on a higher-SES individual’s expectations, while $\beta_2$ tells us how much larger (or smaller) the effect is for lower-SES individuals. I anticipate, based on prior evidence, that $\beta_2$ will be positive when the narrative content is negative (meaning low-SES individuals overreact more to pessimistic news) and possibly also positive for optimistic content if low-SES perhaps also overreact on the upside. In other words, lower-SES groups might exhibit a greater swing in expectations for a given narrative shock. \citet{Polattimur2025} found precisely this pattern for inflation news: less-educated and lower-income households adjusted their inflation expectations more in response to inflation-related news coverage. This project will examine whether the same holds for broader macro narratives and whether it extends to actual decision-making. As a complement, if data permits, I will also investigate partisan heterogeneity (since polarization is inherently political): e.g., do Republicans respond more to Fox narratives and Democrats to CNN/MSNBC narratives, as found in recent studies \citep{Polattimur2025}? While partisan heterogeneity is not the main focus, it will be an important control to ensure that what we attribute to SES is not purely driven by one political group dominating a certain SES category. \section{Expected Contributions and Significance}
The expected findings of this research will have several key implications:
\begin{itemize}
    \item \textbf{Media narratives as a macroeconomic channel:} By demonstrating a link between polarized media narratives and expectation formation, this study will highlight a new mechanism through which media polarization can affect macroeconomic stability. If households in different informational bubbles form divergent expectations, their consumption and investment behaviors may similarly diverge, potentially amplifying regional business cycles or complicating national policy.
    \item \textbf{Evidence for diagnostic expectations:} The project will provide empirical evidence on the relevance of diagnostic expectations in a real-world macro setting. If narrative swings correlate with exaggerated shifts in beliefs that later reverse, it supports the view that representativeness-driven overreactions \citep{Bordalo2022} are at play in household expectations. This bridges the gap between behavioral theory and observed survey data on expectations.
    \item \textbf{Distributional consequences of narratives:} Investigating SES heterogeneity will reveal whether narrative-driven expectation errors have unequal effects across society. Evidence that lower-SES households are more susceptible to biased narratives would suggest that media-driven misperceptions can worsen economic inequality (e.g., through suboptimal financial decisions). This insight could inform targeted communication and education efforts by policymakers to mitigate such disparities.
    \item \textbf{Interdisciplinary methodology:} Methodologically, the project combines text-based narrative analysis with econometric identification in macroeconomics. This interdisciplinary approach can pave the way for future research on the economic impact of information ecosystems and help design interventions (such as public information campaigns) to counteract harmful narrative-induced biases.
\end{itemize}

\bibliographystyle{apalike}
\bibliography{refs}

\end{document}
